Static malware detectors decide whether a file is malicious from its contents alone, before it runs. They come in two
kinds. \emph{Feature-based} detectors parse the file into hand-engineered features, such as the 2{,}568 features of
EMBER~\cite{anderson2018ember, joyce2025ember}, and learn over the feature vector, typically with gradient-boosted
decision trees (GBDTs). \emph{Raw-byte} detectors such as MalConv~\cite{raff2018malware} learn from the bytes
directly. Production engines run both: Microsoft Defender combines models over featurized metadata with deep networks
over full file contents~\cite{microsoft2019insideout}. When a detector misses a malicious file or flags a benign one,
analysts turn to explanations to find out why~\cite{mink2023everybody}.

The two kinds of detector explain themselves in spaces that do not meet. A feature-based detector's explanation
credits features, such as the header checksum or a hashed import bucket, but does not say where in the file that
evidence lies; analysts, who reason about the file itself, must map features back to bytes by hand. A raw-byte
detector's explanation credits bytes, but bytes have no names: an offset means something different in every file, so
these explanations cannot be aggregated across a corpus to reveal what the detector relies on in general. And because
the two kinds of explanation live in different spaces, two detectors cannot be compared on the same file.

Existing work explains each kind of detector in its own space~\cite{fan2020can, warnecke2020evaluating,
coull2019activation, bose2020explaining}. Explaining a feature-based detector directly over the bytes, e.g., by
occluding windows of the file~\cite{zeiler2014visualizing}, costs one feature extraction per window, so byte
resolution is out of reach. Attribution-propagation rules carry explanations through a chain of
models~\cite{shrikumar2017learning, chen2022explaining}, but need every stage to be a model with its own attribution
method, not a parser; FINER~\cite{he2023finer} leaves EMBER-style detectors out, arguing that hashing and aggregation
lose too much to map their explanations back to the program. Attacks that need the mapping from features to bytes
build it by hand, one feature at a time~\cite{pierazzi2020intriguing, severi2021explanation}.

Our key insight is that \emph{every engineered feature is computed from specific bytes, so instrumenting the feature
extractor once records, as a single matrix $P$, which bytes produce each feature and with what weight; and this one
matrix transfers explanations in both directions}. The \emph{pullback} $b = Pa$ places a feature-space explanation on
the bytes, and the \emph{pushforward} $a = P^\top D^\dagger b$ places a byte-level explanation in the features
(Figure~\ref{fig:pullback-pushforward}). We call $P$ the file's \emph{explanation provenance}.

We instrument the EMBER extractor and evaluate both transfers on EMBER-2024 with a GBDT and MalConv, asking the same
three questions of each: is the transferred explanation faithful, does it recover a trigger planted in the training
data, and what security task does it make possible? The pullback is as faithful as occlusion at the same resolution, at
1/135 of its cost, and at byte resolution removes 1.71 times as much evidence; its six top-ranked bytes are exactly a
planted trigger's six bytes in all 50 test files. Used \emph{locally}, it compares the two detectors on a single file:
when the GBDT alone flags malware, MalConv reads a third of the GBDT's evidence the other way, and deleting those bytes
makes MalConv flag 19 of 25 such files, a distinction that vanishes over windows of bytes. The pushforward is faithful
up to the resolution of its source; on MalConv, the features it ranks highest carry about twice as much of the
detector's evidence as those of a surrogate model trained to imitate it. It also turns the planted trigger into a ground
truth for feature-space explainers.
Used \emph{globally}, it gives MalConv's per-file explanations a shared vocabulary, so they can be aggregated over a
corpus: 9 of the 20 fields MalConv relies on most are confirmed, and changing them only as the Windows loader allows
evades 12.4\% of detected held-out malware at a 1\% false-positive rate, twice the rate of as many other loader-ignored
fields. In summary, we make the following contributions:
\begin{itemize}
    \item \textbf{Explanation provenance} (Section~\ref{sec:XML:Framework}): a matrix $P$, recorded by instrumenting a
    feature extractor or estimated by occlusion, and two transfers through it, the pullback and the pushforward.
    \item \textbf{The pullback} (Section~\ref{sec:XML:Pullback}): faithful at a fraction of the cost of direct
    problem-space explanation and precise to the byte; it triages disagreements between a feature-based and a raw-byte
    detector, file by file, where windows of bytes reach the opposite conclusion.
    \item \textbf{The pushforward} (Section~\ref{sec:XML:Pushforward}): faithful up to the resolution of its source,
    and on MalConv more faithful than a surrogate model; it audits a raw-byte detector across a corpus in named features, confirming reliance that, changed only as the
    loader allows, evades an eighth of detected malware; zeroing the audited fields cuts the strongest such attack
    from 16.8\% to 2.9\% at a cost of 3.5 points of detection.
\end{itemize}
